\chapter{Glandular Fever}
\index{Glandular Fever}

\section*{Definition and Overview}
Glandular fever, also known as infectious mononucleosis or "mono", is a common viral illness caused by the Epstein-Barr virus (EBV). It is characterised by a sore throat, fever, swollen lymph glands (hence "glandular" fever), and extreme fatigue. The illness is most common in teenagers and young adults and is often called the "kissing disease" because the virus is spread through saliva. Most people recover fully within a few weeks, but fatigue can persist for weeks to months. There is no specific treatment or vaccine; care focuses on rest, fluids, and symptom relief. The virus stays in the body for life but is usually dormant after the initial illness.

\section*{Epidemiology}
Epstein-Barr virus infects more than 90\% of adults worldwide, usually without causing significant illness in childhood. In many countries and other developed countries, infection is often delayed until adolescence or young adulthood, when it causes symptomatic glandular fever in about 25--50\% of those infected. The condition is most common between the ages of 15 and 25, with a peak around 15--17 years. It affects both sexes equally and occurs throughout the year, with no significant seasonal pattern. Because childhood infection is often silent, most adults have antibodies to EBV by middle age.

\section*{Aetiology and Pathophysiology}
Glandular fever is caused by the Epstein-Barr virus, a member of the herpesvirus family.

\begin{itemize}
    \item \textbf{Transmission:} the virus is spread through saliva, including kissing, sharing drinks, and coughing or sneezing
    \item \textbf{Entry and infection:} the virus infects the lining of the throat and then B-lymphocytes (a type of white blood cell) in the tonsils and lymph nodes
    \item \textbf{Immune response:} the body mounts a strong immune reaction, producing large numbers of atypical lymphocytes and antibodies; this immune response causes the symptoms
    \item \textbf{Latency:} after the initial infection, the virus remains dormant in B-cells for life and can reactivate in saliva intermittently
    \item \textbf{Incubation:} symptoms appear four to six weeks after exposure
    \item \textbf{Infectivity:} the virus is shed in saliva during the illness and for many months afterwards, and intermittently for life
\end{itemize}

\section*{Clinical Features}
Symptoms develop gradually over several days and vary in severity.

\begin{itemize}
    \item Severe sore throat, which can make swallowing painful and is often the most prominent symptom
    \item Fever, often with sweats and chills
    \item Swollen, tender lymph nodes, especially in the neck and armpits
    \item Extreme fatigue and weakness that can last for weeks
    \item Enlarged tonsils, sometimes with a white coating
    \item Headache and body aches
    \item Swollen spleen or liver, and occasionally a faint rash
    \item A rash if amoxicillin or ampicillin antibiotics are taken during the illness
\end{itemize}

\section*{Differential Diagnosis}
Because the symptoms overlap with many other infections, glandular fever must be distinguished from other causes of sore throat and fever.

\begin{itemize}
    \item Streptococcal tonsillitis, a bacterial infection that needs antibiotics
    \item Other viral infections, including cytomegalovirus, HIV seroconversion, adenovirus, and influenza
    \item Toxoplasmosis, a parasitic infection
    \item Viral hepatitis
    \item Acute leukaemia and lymphoma, which can present with fever, lymphadenopathy, and fatigue
    \item Drug reactions and other causes of rash with fever
    \item Chronic fatigue syndrome, if fatigue persists long after the acute illness
\end{itemize}

\section*{When to Seek Medical Care}
See a doctor if you have a severe sore throat, high fever, or swollen glands that do not settle within a few days, or if you feel extremely unwell. People with glandular fever who develop difficulty breathing, severe abdominal pain, or jaundice need urgent review.

\textbf{Contact your local emergency services for:}
\begin{itemize}
    \item Difficulty breathing or swallowing, which may indicate airway obstruction from severe tonsillar swelling
    \item Sudden severe pain in the upper left abdomen, which may indicate a ruptured spleen
    \item Confusion, severe headache, or stiff neck
    \item Jaundice, severe dehydration, or inability to drink fluids
\end{itemize}

\section*{Diagnosis and Investigations}
The diagnosis is usually made from the history and examination, and confirmed with blood tests when needed.

\begin{itemize}
    \item \textbf{Clinical features:} the combination of sore throat, fever, and swollen glands in a young person is suggestive
    \item \textbf{Monospot test:} a rapid blood test for heterophile antibodies, positive in most cases from the second week
    \item \textbf{Full blood count:} shows an elevated white cell count with atypical lymphocytes
    \item \textbf{EBV antibody tests:} specific antibodies distinguish current, recent, and past infection
    \item \textbf{Liver function tests:} mild elevation is common and usually resolves
    \item \textbf{Throat swab:} may be done to exclude strep throat, though the two can coexist
\end{itemize}

\section*{Management}
Treatment is supportive, because antibiotics do not work against the virus.

\begin{itemize}
    \item \textbf{Rest:} rest is important while symptoms are severe, especially early in the illness
    \item \textbf{Pain and fever relief:} paracetamol or ibuprofen for sore throat, fever, and aches
    \item \textbf{Fluids and nutrition:} drink plenty of fluids, and soft, cool foods if swallowing is painful
    \item \textbf{Saltwater gargles and lozenges:} soothe the throat
    \item \textbf{Avoid contact sports:} the spleen can be enlarged and fragile for several weeks, so contact sport is avoided for at least three to four weeks, sometimes longer
    \item \textbf{Gradual return to activity:} resume exercise and work or study gradually as energy returns
    \item \textbf{Corticosteroids:} occasionally used for severe throat swelling or other complications
    \item \textbf{Follow-up:} a doctor may review blood counts or liver function in complicated cases
\end{itemize}

\section*{Complications}
\begin{itemize}
    \item Splenic rupture, a rare but life-threatening emergency usually associated with trauma or strenuous activity
    \item Dehydration from painful swallowing
    \item Airway obstruction from massively enlarged tonsils
    \item Secondary bacterial throat infection
    \item Liver inflammation (hepatitis) with jaundice
    \item Anaemia and low platelets in rare cases
    \item Neurological complications, including meningitis, encephalitis, and Guillain-Barré syndrome, which are rare
    \item Prolonged fatigue, which can persist for weeks or months
    \item A rash if ampicillin-type antibiotics are taken during the illness
\end{itemize}

\section*{Prognosis and Outcomes}
Most people with glandular fever recover fully within two to four weeks, though fatigue commonly lingers for several more weeks and occasionally for months. The illness is almost never fatal in healthy people, and the virus remains dormant for life without causing further illness in the vast majority. People can return to normal activity gradually as their energy returns, with the important caveat of avoiding contact sport until the spleen has recovered. Serious complications are rare, and immunity after the illness is lifelong.

\section*{Prevention and Self-Care}
\begin{itemize}
    \item There is no vaccine or medicine that prevents EBV infection, and most adults have already been infected
    \item Avoid sharing drinks, eating utensils, and toothbrushes with someone who is unwell
    \item Practise good hand hygiene, especially after contact with saliva
    \item Rest and drink plenty of fluids during the illness
    \item Take paracetamol or ibuprofen for pain and fever as directed
    \item Avoid strenuous exercise and contact sport for at least three to four weeks after the illness
    \item Do not take antibiotics unless prescribed, and tell the doctor about the glandular fever if antibiotics are being considered
    \item Return to work, study, and exercise gradually as your energy improves
\end{itemize}

\section*{Sources and Further Reading}
The following reputable sources provide further information for healthcare students and patients seeking reliable, current advice about glandular fever.

\begin{itemize}
    \item Healthdirect Australia. \textit{Glandular fever.} https://www.healthdirect.gov.au/
    \item Australian Government Department of Health and Aged Care. \textit{Glandular fever (infectious mononucleosis).} https://www.health.gov.au/
    \item NSW Health. \textit{Glandular fever fact sheet.} https://www.health.nsw.gov.au/
    \item NHS. \textit{Glandular fever.} https://www.nhs.uk/
    \item Centers for Disease Control and Prevention. \textit{Epstein-Barr virus and infectious mononucleosis.} https://www.cdc.gov/
    \item Mayo Clinic. \textit{Infectious mononucleosis.} https://www.mayoclinic.org/
\end{itemize}
